\documentclass[11pt,a4paper]{article}

\usepackage[T1]{fontenc}
\usepackage[utf8]{inputenc}
\usepackage{lmodern}
\usepackage[a4paper,margin=26mm]{geometry}
\usepackage{amsmath,amssymb,amsthm,mathtools}
\usepackage{booktabs,array}
\usepackage{microtype}
\usepackage{xcolor}
\usepackage{enumitem}
\usepackage{url}
\usepackage[colorlinks=true,linkcolor=blue!55!black,citecolor=blue!55!black,
            urlcolor=blue!55!black]{hyperref}

\newtheorem{theorem}{Theorem}
\newtheorem{lemma}[theorem]{Lemma}
\newtheorem{corollary}[theorem]{Corollary}
\newtheorem{remark}[theorem]{Remark}
\newcommand{\Z}{\mathbb{Z}}
\newcommand{\Sign}{\mathsf{Sign}}
\newcommand{\Verify}{\mathsf{Verify}}
\newcommand{\pk}{\mathsf{pk}}
\newcommand{\sk}{\mathsf{sk}}
\newcommand{\bits}{\operatorname{bits}}
\newcommand{\Hash}{\mathcal{H}}

\title{\bfseries Response-Rescaling Fresh-Message Forgeries\\
for SQIsignTriangle\\[0.3em]
\large A Conditional One-Signature Transformation Across All Four Submitted
Parameter Sets}
\author{Martin Feussner\\
\small Selmer Center, University of Bergen\\
\small \texttt{martin.feussner@uib.no}}
\date{27 September 2026}

\begin{document}
\maketitle

\begin{center}
\fbox{\parbox{0.92\linewidth}{\small\textbf{Provenance and review status.}
The attack was independently found by an AI system, \mbox{OpenAI Codex} using
Daybreak Blue at maximum reasoning effort, without a human first proposing
this specific attack. The AI produced the derivation, implementation,
experiments, and initial manuscript. At the date above, no human
cryptanalyst has independently verified the argument, implementation,
measurements, or novelty assessment. This preliminary disclosure is provided
for human review and reproduction.}}
\end{center}
\vspace{0.7em}

\begin{abstract}
SQIsignTriangle represents a response by an auxiliary curve, an auxiliary
$2^e$-torsion basis, and an odd integer $q$. Verification reconstructs a
two-dimensional kernel generated by
\[
  ([q]P_{\rm pk},P_{\rm aux})
  \quad\text{and}\quad
  ([q]Q_{\rm pk},Q_{\rm aux}),
\]
recovers two possible commitment curves, and accepts when $q\equiv c_2
\pmod{c_1}$ for the hash challenge of either curve and the message.

This response has an unchecked scaling symmetry. For any odd unit $u$ modulo
$2^e$, replacing $q$ by $q'\equiv uq\pmod{2^e}$ and multiplying the auxiliary
basis by $u$ multiplies both kernel generators by the same unit. The same
kernel gives an isomorphic quotient whose two codomain factors have the same
unordered isomorphism classes, and hence the same $j$-invariants. Starting
from one valid signature, an attacker can select an odd $e$-bit integer $q'$
satisfying a fresh-message hash congruence, set
$u=q'q^{-1}\bmod 2^e$, and scale the four serialized auxiliary-basis
coefficients. This one-signature transformation is proved when the selected
target challenge satisfies the public sufficient condition
$2^{e-1}>2c_1$. It uses one signing query and requires neither secret
information nor a hash weakness.

We reproduced the attack against the unmodified submitted reference verifier
on three independently generated keys for each of the 128-, 160-, 256-, and
512-bit parameter sets. The sufficient condition held in every trial, one
signing query was used per trial, and all 12 fresh-message forgeries were
accepted. The slowest forgery computation took 4.33 seconds on a four-core
Intel Core i5-6500, and the other three levels took below 0.7 seconds. The
submission provides no lower bound or lower-tail distribution for honest
response lengths, so we do not claim that the condition is guaranteed or
holds with overwhelming probability. If it fails, the general attack repeats
signing queries until it obtains an eligible response; this report does not
bound the expected number of queries.
\end{abstract}

\section{Introduction}

SQIsignTriangle is an isogeny-based signature submission to the
Next-generation Commercial Cryptographic Algorithms Program (NGCC). It follows a
Fiat--Shamir transformation of a three-move protocol and replaces the usual
fixed challenge isogeny with a modular degree condition
\begin{equation}
                         q\equiv c_2\pmod{c_1}.
  \label{eq:challenge}
\end{equation}
The response is represented through Kani's lemma by a
$2^e$-isogeny from a product $E_{\rm pk}\times E_{\rm aux}$, where
$e=\bits(q)$. The verifier reconstructs its kernel from the public basis, the
integer $q$, and an auxiliary basis contained in the signature
\cite{sqitriangle-spec}.

The kernel representation is homogeneous: multiplying every kernel
generator by a unit modulo $2^e$ does not change the generated subgroup. In
the submitted response format, an attacker can perform exactly this scaling
while changing the integer used in Equation~\eqref{eq:challenge}. The
geometric part of verification continues to recover the original commitment,
whereas the replacement integer can be chosen to satisfy a fresh-message
challenge whenever the public response-size condition below holds.

The attack is especially direct:
\begin{itemize}[leftmargin=1.5em]
  \item it uses one ordinary signature whenever the public response-size
        condition holds, and no secret-key information;
  \item it targets the specified verification equation and succeeds without
        malformed inputs, undefined behavior, or a weak hash function;
  \item its only expensive step is one public verifier-style isogeny-chain
        computation; and
  \item the same transformation applies to all four submitted parameter sets.
\end{itemize}

This is stronger than ordinary signature malleability. The output verifies on
a message that was never submitted to the signing oracle. In every reported
control, the genuine signature rejected the target message and the transformed
signature rejected the source message. Those two rejection observations are
experimental controls, not consequences required by the algebraic theorem.

\paragraph{Related work and novelty.}
The NGCC tracker credits Tako Boris Fouotsa with two earlier SQIsignTriangle
findings \cite{ngcc-tracker}. Finding \texttt{sign-27-3} is already a
fresh-message forgery: it keeps a known valid response fixed and searches over
candidate fresh messages until the resulting challenge satisfies
$q\equiv c_2\pmod{c_1}$. Because the relevant challenge value has
$\lambda/2$ bits, this requires approximately $2^{\lambda/2}$ classical hash
trials. Finding \texttt{sign-27-4} observes that distinct modular challenges
can share an identical response, invalidating the special-soundness argument.

The present attack is therefore not the first forgery result for
SQIsignTriangle. It exploits a different symmetry in the same modular-challenge
weakness. Instead of fixing the response and searching for a compatible
message, it fixes an arbitrary fresh message, changes $q$ to meet that
message's challenge, and rescales the auxiliary basis so that the verification
kernel is preserved. Subject to the public response-size condition, this
removes the $2^{\lambda/2}$ hash search and reduces the work to a
verifier-style computation plus modular arithmetic. We characterize the result
as a substantially faster, practically demonstrated exploitation and a
practical strengthening of the same modular-challenge weakness, independently
derived through response rescaling: all four submitted parameter sets were
forged within seconds in our experiments. A public discussion of malleability
in standard SQIsign concerns strong unforgeability rather than this
fresh-message mechanism \cite{sqisign-suf}.

\section{The verification equation}

We state only the parts of SQIsignTriangle needed for the attack. Let the
field parameter satisfy $p=c\mathbin{\cdot}2^f-1$. A public key contains a curve
$E_{\rm pk}$ and a hint from which the verifier reconstructs a basis
$(P^{(f)}_{\rm pk},Q^{(f)}_{\rm pk})$ of the full $2^f$-torsion. A signature is
\begin{equation}
           \sigma=(E_{\rm aux},P_{\rm aux},Q_{\rm aux},q),
  \label{eq:sig}
\end{equation}
where $q$ is positive and odd and
\begin{equation}
                         e=\bits(q).
  \label{eq:e}
\end{equation}
The submitted verifier also checks $3\le e\le f$ and verifies that the public
and auxiliary pairs have exact order $2^e$.

Define the cofactor-reduced public basis
\[
  \bar P_{\rm pk}=[2^{f-e}]P^{(f)}_{\rm pk},\qquad
  \bar Q_{\rm pk}=[2^{f-e}]Q^{(f)}_{\rm pk}.
\]
Scalar multiplication commutes with the cofactor reduction used in the
implementation. The verifier constructs the subgroup
\begin{equation}
 K_\sigma=\left\langle
   ([q]\bar P_{\rm pk},P_{\rm aux}),
   ([q]\bar Q_{\rm pk},Q_{\rm aux})
 \right\rangle
 \subseteq
 (E_{\rm pk}\times E_{\rm aux})[2^e].
 \label{eq:kernel}
\end{equation}
It evaluates the corresponding two-dimensional isogeny chain and obtains a
split codomain $F_1\times F_2$. For $i\in\{1,2\}$ it derives
\begin{equation}
 (c_{1,i},c_{2,i})=
 \Hash(j(E_{\rm pk}),j(F_i),M).
 \label{eq:hash}
\end{equation}
It accepts if, for either component,
\begin{equation}
                         q\equiv c_{2,i}\pmod{c_{1,i}}.
 \label{eq:accept}
\end{equation}
The submitted implementation tries $F_1$ first and then $F_2$. Thus no
component label needs to be predicted by the attacker.

There is a presentation inconsistency in Algorithm~4.3 of the specification:
its pseudocode hashes an as-yet undefined $E_{\rm com}$ before reconstructing
$F_1\times F_2$. Section~4.4 nevertheless states the intended relation: the
kernel in Equation~\eqref{eq:kernel} is evaluated, one recovered factor is the
commitment, and the response must meet the modular challenge. The reference
verifier resolves the ordering issue by reconstructing the quotient first and
then evaluating Equation~\eqref{eq:accept} for both factors. Our algebraic
argument targets this intended specified relation. Claims about byte decoding,
branch order, and measured acceptance refer separately to the submitted
implementation.

\subsection{Serialized auxiliary bases}

For each submitted level, the signature stores $E_{\rm aux}$, four little-endian
coefficients $(r_1,r_2,s_1,s_2)$, and $q$. Let
$(P^{(f)}_{\rm aux},Q^{(f)}_{\rm aux})$ be the deterministic full basis used by
the decoder. Up to the implementation's equivalent order of operations, the
decoded points are
\begin{align}
 P_{\rm aux}&=[2^{f-e}]
   ([r_1]P^{(f)}_{\rm aux}+[r_2]Q^{(f)}_{\rm aux}),\label{eq:decode-p}\\
 Q_{\rm aux}&=[2^{f-e}]
   ([s_1]P^{(f)}_{\rm aux}+[s_2]Q^{(f)}_{\rm aux}).\label{eq:decode-q}
\end{align}
Consequently, multiplying all four coefficients by the same value modulo
$2^e$ multiplies both decoded points by that value. The differences introduced
by reduction modulo $2^e$ vanish after multiplication by $2^{f-e}$.

The submitted byte lengths are shown in Table~\ref{tab:params}. They agree
between the technical specification and the reference API.

\begin{table}[t]
\centering
\caption{Submitted SQIsignTriangle parameter and encoding sizes.}
\label{tab:params}
\begin{tabular}{@{}lrrrrr@{}}
\toprule
Instance & $\lambda$ & $f$ & Public key & Signature & Coefficient/$q$ \\
         &           &     & bytes      & bytes     & bytes \\
\midrule
lvl1 & 128 & 248  & 65  & 204 & 28  \\
lvl2 & 160 & 309  & 81  & 255 & 35  \\
lvl5 & 256 & 500  & 129 & 408 & 56  \\
lvl6 & 512 & 1016 & 257 & 816 & 112 \\
\bottomrule
\end{tabular}
\end{table}

\section{The response-rescaling symmetry}

Write $R_e=\Z/2^e\Z$. Since the verifier requires $q$ to be odd, $q$ is a
unit in $R_e$.

\begin{lemma}[Kernel invariance]
\label{lem:kernel}
Let $q'$ be any odd integer with $\bits(q')=e$, and set
\begin{equation}
                     u=q'q^{-1}\pmod{2^e}.
 \label{eq:u}
\end{equation}
Replace the auxiliary basis by
\[
  (P'_{\rm aux},Q'_{\rm aux})=([u]P_{\rm aux},[u]Q_{\rm aux}).
\]
Then the kernel reconstructed from $(q',P'_{\rm aux},Q'_{\rm aux})$ is exactly
$K_\sigma$.
\end{lemma}

\begin{proof}
Because $q'\equiv uq\pmod{2^e}$ and all points in
Equation~\eqref{eq:kernel} are in the $2^e$-torsion, the two new generators are
\begin{align*}
 ([q']\bar P_{\rm pk},P'_{\rm aux})
   &=([uq]\bar P_{\rm pk},[u]P_{\rm aux})
     =[u]([q]\bar P_{\rm pk},P_{\rm aux}),\\
 ([q']\bar Q_{\rm pk},Q'_{\rm aux})
   &=([uq]\bar Q_{\rm pk},[u]Q_{\rm aux})
     =[u]([q]\bar Q_{\rm pk},Q_{\rm aux}).
\end{align*}
Both $q$ and $q'$ are odd, so $u$ is odd and hence a unit in $R_e$.
Multiplication by $u$ is an automorphism of the $2^e$-torsion. The scaled
generators therefore span the same subgroup. The exact-order checks also
continue to hold.
\end{proof}

\begin{corollary}[Quotient and factor invariance]
The original and rescaled signatures have the same kernel and therefore define
isomorphic quotients. Whenever the quotient is represented as a split product,
the two codomain factors have the same unordered isomorphism classes. In
particular, the unordered pair of $j$-invariants used in
Equation~\eqref{eq:hash} is unchanged.
\end{corollary}

\begin{proof}
The source product is unchanged and Lemma~\ref{lem:kernel} gives equality of
the finite kernel subgroups. The quotient by that subgroup is consequently
unique up to isomorphism. Its decomposition yields the same factor isomorphism
classes up to ordering, so their $j$-invariants agree.
\end{proof}

This is a property of the mathematical response representation. It does not
depend on lenient parsing or on a particular sequence of formulas inside the
reference code.

\section{Fresh-message forgery}

Let $M^*$ be any message distinct from the signing query. The attacker first
runs the public part of verification on a valid signature to recover
$F_1,F_2$. Select a component $F_i$ and compute the target challenge
$(c_1,c_2)$ from Equation~\eqref{eq:hash} using $M^*$.

The attacker needs an odd integer in the interval
\begin{equation}
       2^{e-1}\le q'<2^e,\qquad q'\equiv c_2\pmod{c_1}.
 \label{eq:qprime}
\end{equation}
The challenge modulus $c_1$ is an odd prime. A direct construction is
\begin{align}
 k_0&=\left\lceil\frac{2^{e-1}-c_2}{c_1}\right\rceil,\\
 q_0&=c_2+k_0c_1.
\end{align}
Use $q'=q_0$ if it is odd and $q'=q_0+c_1$ otherwise. The exact public
eligibility test is whether one of these odd representatives lies in the
$e$-bit interval. A simple sufficient condition is
\begin{equation}
                         2^{e-1}>2c_1
 \label{eq:room}
\end{equation}
because it leaves room for two consecutive representatives. In the submitted
sets, $c_1$ has at most
$64,80,128,$ and $256$ bits, respectively. Our genuine signatures had
$e$ between $219$ and $895$, depending on the level, so the number of
available representatives is enormous.

Equation~\eqref{eq:room} is not guaranteed by honest signing in either the
specification or the submitted code. The specification's discussion of
\texttt{EquivalentSpecialIdeal} gives heuristic typical-size reasoning and an
upper-bound argument for its lattice output, but no lower bound or lower-tail
probability for $q$. The signer computes $e=\bits(q)$ and rejects responses
that are even or exceed the torsion and encoding limits; it does not reject a
response merely because $e$ is small relative to $c_1$. Consequently, the
available material does not justify saying that
Equation~\eqref{eq:room} holds with overwhelming probability. Its satisfaction
in the experiments below is evidence about those tested executions only.

After finding $q'$, compute $u$ by Equation~\eqref{eq:u}. Keep
$E_{\rm aux}$ unchanged, replace the four serialized basis coefficients by
\begin{equation}
 (r_1',r_2',s_1',s_2')
 =u(r_1,r_2,s_1,s_2)\pmod{2^e},
 \label{eq:coeff}
\end{equation}
and replace $q$ by $q'$. Because $q'$ has the same bit length, the decoder
uses the same cofactor and all transformed coefficients fit the submitted
encoding.

\begin{theorem}[Conditional one-signature forgery]
\label{thm:forgery}
Suppose a valid SQIsignTriangle signature has response bit length $e$ and the
challenge of a selected recovered factor on $M^*$ satisfies
Equation~\eqref{eq:room}. The transformation in
Equations~\eqref{eq:qprime}--\eqref{eq:coeff} produces a valid signature on
the chosen target message $M^*$ under the verification relation in
Equations~\eqref{eq:kernel}--\eqref{eq:accept}.
\end{theorem}

\begin{proof}
The constructed $q'$ is positive, odd, has bit length $e$, and satisfies the
target congruence. Lemma~\ref{lem:kernel} shows that the verifier reconstructs
the same kernel. It therefore obtains the same codomain factors, including
the selected $F_i$. When it hashes $F_i$ with $M^*$, it obtains the public
values $(c_1,c_2)$ used to construct $q'$, so
$q'\equiv c_2\pmod{c_1}$. All range, curve, basis-order, isogeny-chain, and
congruence checks pass.
\end{proof}

The theorem permits either recovered factor to be selected, provided its own
target challenge meets Equation~\eqref{eq:room}. This is algebraic rather than
an inference from the experiments: quotient invariance preserves both factor
$j$-invariants, and Equation~\eqref{eq:accept} is a disjunction over the two
branches. The separate tests below exercise this conclusion in the submitted
implementation.

\paragraph{Attack algorithm.}
The complete repeat-until-eligible adversary is as follows.
\begin{enumerate}[leftmargin=2em]
  \item Choose a target message $M^*$ and request a signature $\sigma$ on a
        distinct source message.
  \item Parse $q$, compute $e=\bits(q)$, and execute the public
        isogeny-chain portion of verification to obtain $F_1,F_2$.
  \item Hash each recovered factor with $M^*$. If neither branch has an odd
        $e$-bit representative in its target residue class, discard this
        response and repeat from Step~1 with another source signing query.
        Equation~\eqref{eq:room} is a sufficient public test for a branch.
  \item Construct $q'$ by Equation~\eqref{eq:qprime}, compute
        $u=q'q^{-1}\bmod 2^e$, and apply Equation~\eqref{eq:coeff}.
  \item Output the unchanged auxiliary curve, the four transformed
        coefficients, and $q'$.
\end{enumerate}
For an eligible first response, this uses exactly one signing query. Otherwise
it makes further queries. Because the honest-response lower tail is not
specified or proved, we give no theorem for the expected query count. The
source message can also be used to identify an originally accepting factor by
testing Equation~\eqref{eq:accept}, although the submitted verifier's two
target branches make that identification unnecessary.

\subsection{EUF-CMA interpretation}

In the EUF-CMA experiment, the adversary may request signatures on messages of
its choice and wins by producing a valid signature on a message absent from
its query set \cite{gmr}. The algorithm above outputs a signature on the
unqueried $M^*$ and is therefore a fresh-message forgery rather than a second
encoding for a queried message. Theorem~\ref{thm:forgery} proves one-query
success only when the first response meets the stated public condition. A
formal asymptotic EUF-CMA advantage for the repeat-until-eligible adversary
would additionally require a lower bound on the probability of that event;
the submission does not provide one, and this paper does not infer one from
12 trials. The attack does not recover the secret key.

\section{Experimental validation}

We built the portable reference implementation from the archived submission
and called its unmodified signature-verification API. The experiment used GCC
13.3.0 on Ubuntu 24.04, with an Intel Core i5-6500 at 3.20 GHz and four physical
cores. Only one core was needed by each attack run.

For every parameter set we generated three fresh key pairs. Before each key
generation, the harness DRNG was initialized with a distinct 48-byte seed
obtained as
\[
 \operatorname{SHAKE256}_{48}(\texttt{SQIsignTriangle level L, key I}),
\]
where $L\in\{1,2,5,6\}$ and $I\in\{0,1,2\}$. For each key we signed the literal
source message
\[
 \texttt{SQIsignTriangle source, level L, key I}
\]
and forged the distinct target
\[
 \texttt{SQIsignTriangle fresh target, level L, key I}.
\]

In each trial we observed all four of the following controls:
\begin{enumerate}[leftmargin=2em]
  \item the genuine signature was accepted on the source message;
  \item the genuine signature was rejected on the target message;
  \item the forged signature was accepted on the target message; and
  \item the forged signature was rejected on the source message.
\end{enumerate}
The theorem requires the third observation under its stated condition. It
does not require the second or fourth: a congruence coincidence could make a
genuine signature accept another message or a transformed signature accept
the source. Those two rejections report what the reference verifier did in
these experiments. All 12 trials passed. Table~\ref{tab:results} reports the observed response
bit lengths and the time for public codomain reconstruction, target challenge
derivation, modular rescaling, and serialization. It excludes the signing
oracle's work and a final verification call.

\begin{table}[t]
\centering
\caption{End-to-end fresh-message forgery results on three independent keys
per submitted level.}
\label{tab:results}
\begin{tabular}{@{}lrrrr@{}}
\toprule
Instance & Accepted & Observed $e$ & Forgery time range & Maximum \\
\midrule
lvl1 & 3/3 & 219--221 & 0.033--0.046 s & 0.046 s \\
lvl2 & 3/3 & 275--276 & 0.023--0.051 s & 0.051 s \\
lvl5 & 3/3 & 442--443 & 0.233--0.690 s & 0.690 s \\
lvl6 & 3/3 & 893--895 & 0.401--4.329 s & 4.329 s \\
\bottomrule
\end{tabular}
\end{table}

As a separate component-branch check, on one key at each level we deliberately
constructed a second forgery using the codomain factor that did not satisfy
the original source-message congruence. All four of those signatures were
also accepted on their target messages. This experimentally exercises both
branches of the submitted verifier on those four instances. The general
factor-choice statement used by the theorem follows separately from quotient
invariance and the disjunctive relation in Equation~\eqref{eq:accept}, with
the response-size condition applied to the chosen branch.

For a deterministic level-1 check, we also seeded the DRNG with the 48 bytes
\texttt{00,01,...,2f} and signed
\texttt{SQIsignTriangle source message}. The signature had $e=220$. For the
target \texttt{SQIsignTriangle fresh target message}, the second recovered
factor gave
\[
 c_1=18337881540236675959,
 \qquad c_2=5510551482462846558.
\]
The construction produced an odd 220-bit $q'$ in that residue class. In this
test, the reference API accepted the transformed 204-byte signature on the
target and rejected it on the source; it accepted the genuine signature on
the source and rejected it on the target.

\section{Root cause and scope}

The response is checked in two logically separate ways:
\begin{enumerate}[leftmargin=2em]
  \item the kernel determines a quotient whose factor must reproduce the
        commitment curve; and
  \item the integer $q$ must satisfy a short modular challenge relation.
\end{enumerate}
The verifier assumes these checks refer to one fixed response isogeny. In
fact, the kernel is invariant under simultaneous multiplication of its public
and auxiliary coordinates by a unit, while the integer representative $q$
changes. This lets the attacker retain the commitment while selecting the
integer representative after seeing the target challenge.

The attack is independent of the submitted hash implementation: it only
evaluates the public hash and uses its output. It also survives strict and
canonical byte decoding, because every transformed coefficient and $q'$ is
within its prescribed range. The previously reported assertion failure on an
all-zero signature is unrelated \cite{ngcc-tracker}.

The proof of EUF-CMA security must account for the full unit orbit of a
response. Treating $q$ as the norm selected by the honest
\texttt{EquivalentSpecialIdeal} algorithm is insufficient for adversarial
verification, because the verifier does not establish that provenance. The
same symmetry follows from the kernel and modular challenge stated in
Section~4.4 and Algorithm~4.3 of the specification. We therefore describe the
point-level vulnerability as design-level with respect to the intended
verification relation. Algorithm~4.3 itself has the ordering ambiguity noted
above; the claims that the verifier hashes both factors in sequence and that
coefficient rescaling realizes point rescaling are supported by the submitted
implementation and encoding description. The experiments establish that the
reference code implements the vulnerable relation; the argument does not rely
on a wrapper defect, malformed parsing, or undefined behavior.

\paragraph{Repair direction.}
A repair must cryptographically bind the scalar representative used in the
challenge check to the geometric response. Merely checking point orders,
requiring canonical byte strings, or changing the hash function does not
remove the unit orbit. Possible redesigns include a publicly verifiable
canonical representative for the entire response orbit, a proof that binds
the claimed norm to the represented response, or a challenge that fixes the
full scalar rather than only its residue modulo a much smaller prime. Each
option changes the protocol and requires a new correctness and security
analysis; adding $q$ to the existing pre-response hash would introduce a
circular dependency and is not by itself a complete fix.

\section{Conclusion}

SQIsignTriangle's compact two-dimensional response admits a public unit
rescaling that leaves its verification kernel unchanged. The modular-degree
challenge is applied to a scalar representative rather than to the invariant
kernel, so an eligible representative can be retargeted to a fresh-message
hash challenge. One genuine signature was sufficient in every reported trial,
the work was comparable to a verification, and the attack was reproduced on
all four submitted parameter sets. The proof of one-query success is
conditional on Equation~\eqref{eq:room}; neither a universal guarantee nor an
overwhelming-probability bound for honest responses is claimed. The result
exhibits concrete fresh-message forgeries and a design-level verification
symmetry, while a fully quantified asymptotic EUF-CMA advantage would require
an honest-response eligibility bound absent from the submission.

\begin{thebibliography}{9}

\bibitem{sqitriangle-spec}
Wei Yu, Bin Wang, Tingyu Xing, Wenbiao Chen, Jiaxuan Qu, Xiuyu Qiu,
Song Tian, and Kunpeng Wang.
\newblock \emph{SQIsignTriangle: Short Quaternion and Isogeny Signature with a
Triangle-shaped Framework}.
\newblock NGCC submission specification, 2026.
\newblock \url{https://www.niccs.org.cn/niccs/Proposal/Public-Key%20Cryptographic%20Algorithms/Round%201%20candidates/SQIsignTriangle.zip}.

\bibitem{ngcc-tracker}
NGCC Security Project.
\newblock \emph{SQIsignTriangle (sign-27)}.
\newblock Public issue tracker, including findings \texttt{sign-27-3} and
\texttt{sign-27-4} credited there to Tako Boris Fouotsa; accessed 27 September
2026.
\newblock \url{https://ngcc.dev/reports/sign-27.html}.

\bibitem{sqisign-suf}
Dustin Ray.
\newblock \emph{Towards SUF-CMA SQIsign}.
\newblock Public issue discussion, June 2026.
\newblock \url{https://github.com/anchorageoss/sqisign-rs/issues/19}.

\bibitem{fs}
Amos Fiat and Adi Shamir.
\newblock How to prove yourself: Practical solutions to identification and
signature problems.
\newblock In \emph{Advances in Cryptology -- CRYPTO '86}, pages 186--194.
Springer, 1987.

\bibitem{gmr}
Shafi Goldwasser, Silvio Micali, and Ronald L. Rivest.
\newblock A digital signature scheme secure against adaptive chosen-message
attacks.
\newblock \emph{SIAM Journal on Computing}, 17(2):281--308, 1988.

\end{thebibliography}

\end{document}
